\documentclass[a4paper,12pt]{article}
\usepackage[utf8]{inputenc}
\usepackage{amsmath, amssymb}
\usepackage{geometry}
\usepackage{graphicx}
\usepackage[utf8]{inputenc}
\usepackage[T1]{fontenc}
\usepackage[spanish]{babel}
\usepackage{url}
\geometry{margin=1in}

% Datos de la tarea
\title{\textbf{Seminario Matemáticas Aplicadas II:\\Consecuencias estadísticas de las distribuciones de colas gordas en la ciencia y la toma de decisiones}}
\author{Eliud Daniel Santillan Aparicio 415069696 \\ Leonardo Vigueras Salazar 310738244}
%\date{\today}
\date{}


\begin{document}

\maketitle
\section*{Instrucciones}

\subsection*{Instrucciones para Diseñar y Resolver un ABP Usando ChatGPT}

\textbf{Aprendizaje Basado en Problemas (ABP)}, aprovechando a ChatGPT como herramienta principal. Este ejercicio sigue tres fases: diseño del ABP, resolución del ABP usando el esquema de Polya y la entrega de los resultados.

\section*{Fase 1: Diseña un ABP Personalizado}

\begin{enumerate}
    \item \textbf{Elige una Tarea:} Case study: \textit{Estrategias antifrágiles en el contexto de monopolios.}
    \item \textbf{Contexto y los Objetivos:}
    \begin{itemize}
    \item \textbf{Contexto:} ¿Cómo crear un monopolio? ¿Los monopolios se crean o surgen? Es decir, ¿son impuestos por una voluntad individual suficientemente poderosa como para imponerlos (p.\,ej., un Estado) o emergen de forma orgánica bajo ciertas condiciones de oferta, demanda y ausencia de regulación?

    Las hipótesis iniciales de este proyecto postulan que los sistemas económicos son comparables a los ecosistemas biológicos y, por lo tanto, susceptibles de análisis mediante las mismas metodologías de sistemas dinámicos que habitualmente se emplean para su estudio y modelación.

    Por ello, exploraremos hasta qué punto es posible aislar y modelar el fenómeno de los monopolios como resultado de condiciones sistémicas, y determinaremos si, en este contexto, existe antifragilidad.
    \item \textbf{Objetivo:} Proponer un modelo que formalice las condiciones sistémicas que permiten (y tal vez inducen) dinámicas monopólicas, aislando e identificando monopolios posibles y existentes; entrenarlo y validarlo con ejemplos y datos reales, e investigar y concluir si en dicho modelo existen características antifrágiles.
    \end{itemize}
    \item \textbf{Diseña Prompts Iniciales:}
    Se listan los prompts pero su desarrollo se anexa en un documento aparte llamado: promptsIniciales.pdf
    \begin{enumerate}
            \item \textbf{Exploración conceptual y teórica}
    \begin{itemize}
        \item Analizar, con ejemplos históricos, si los monopolios han surgido más comúnmente por imposición (top-down) o de forma emergente (bottom-up), y cuáles han sido las condiciones que facilitaron cada tipo.
        \item Establecer paralelismos entre dinámicas ecológicas—como la selección natural, la competencia o el mutualismo—y las dinámicas económicas en mercados no regulados.
        \item Identificar características de ecosistemas biológicos que puedan compararse con mercados económicos para estudiar la emergencia y consolidación de monopolios.
        \item Revisar autores y marcos teóricos que hayan comparado explícitamente sistemas económicos con sistemas biológicos o complejos, y sintetizar los modelos o resultados propuestos.
    \end{itemize}

    \item \textbf{Diseño del modelo sistémico}
    \begin{itemize}
        \item Proponer un esquema para diseñar un modelo dinámico—ya sea mediante ecuaciones diferenciales o autómatas celulares—que represente competencia entre agentes económicos con posibilidad de aparición de monopolios.
        \item Determinar las variables sistémicas clave para simular la evolución de un mercado hacia una concentración monopólica (por ejemplo: costos marginales, elasticidad de demanda, acceso a información, etc.).
        \item Identificar datasets, tanto reales como simulados, que puedan alimentar un modelo orientado a detectar condiciones inductoras de monopolios.
        \item Formalizar computacionalmente el concepto de “emergencia orgánica” de un monopolio en ausencia de intervención estatal.
    \end{itemize}

    \item \textbf{Validación empírica}
    \begin{itemize}
        \item Establecer criterios e indicadores históricos o financieros para identificar empíricamente si un monopolio emergió por condiciones sistémicas o por intervención estatal.
        \item Presentar ejemplos de modelos validados en la literatura científica que buscan predecir concentración de mercado o formación de monopolios.
        \item Seleccionar métricas que permitan medir dinámicamente la evolución hacia estructuras de mercado monopólicas, e investigar analogías con indicadores como entropía o biodiversidad en ecología.
    \end{itemize}

    \item \textbf{Evaluación de antifragilidad}
    \begin{itemize}
        \item Definir y operacionalizar el concepto de antifragilidad dentro de un modelo económico que simula la formación de monopolios.
        \item Explorar si un sistema que tiende a formar monopolios puede considerarse antifrágil, y bajo qué condiciones se diferencian la antifragilidad del sistema de la del agente que adquiere poder monopólico.
        \item Diseñar pruebas de estrés para un modelo de formación de monopolios con el fin de evaluar si su comportamiento ante perturbaciones es antifrágil o simplemente robusto.
        \item Estudiar analogías entre especies invasoras en ecología y actores económicos que alcanzan posiciones monopólicas, y extraer lecciones relevantes sobre antifragilidad desde esa perspectiva.
    \end{itemize}

    \end{enumerate}
    \item \textbf{Crea el Escenario del ABP:}
    \paragraph{Describe detalladamente el problema, asegurándote de incluir los datos, contexto y roles si aplica:}
    \begin{itemize}
        \item Diseñar un framework formal de naturaleza matemática que establezca un sistema coherente de definiciones, objetos, relaciones y funciones, con el objetivo de sustentar el entrenamiento de un modelo capaz de identificar si, dado un conjunto de información suficientemente representativo, existe en el sistema modelado un espacio propicio para la emergencia de un monopolio.

Idealmente, se busca construir una representación del monopolio como un objeto matemático bien definido dentro del sistema dinámico considerado. Intuitivamente, el propósito es alcanzar una formulación que permita enunciar afirmaciones del tipo:
\begin{quote}
    Para un sistema económico \( X \), parametrizado por el vector \( \vec{Y} \), tras \( n \) iteraciones bajo ciertas reglas de evolución, emerge un objeto \( \mathcal{M} \) que exhibe propiedades compatibles con la noción formal de monopolio.
\end{quote}


    \end{itemize}

    \item \textbf{Planifica las Tareas Específicas:} referir a documento anexo "Planeación.pdf".

\end{enumerate}




\section*{Fase 2: Resuelve el ABP Usando el Esquema de Polya}

Aplica los pasos del esquema de resolución de Polya para resolver el problema. Documenta cada etapa y utiliza ChatGPT como apoyo. El proceso debe estar documentado. 


\subsubsection*{1. Plan (Desarrollado en Fase 1)}
    \textbf{Tareas:}
\begin{itemize}
    \item \textbf{Etapa I: Fundamentación Teórica y Conceptual} Realizar una revisión bibliográfica sistemática sobre:
    \begin{itemize}
        \item \textbf{Modelos de formación de monopolios:} Económicos, históricos y regulatorios: Después de una exhaustiva investigación (la bibliografía usada queda anexa en el repositorio, en el documento: "bilbiografiaadp1.pdf") se llego a la conclusion de que a lo largo de la historia han existido varios tipos de monopolios, cada uno definido por cómo se logra el control exclusivo de la disposición del recurso en cuestión al mercado:
\begin{itemize}
    \item \textbf{Monopolio natural:} Se presenta cuando un único proveedor resulta ser el más eficiente debido a altos costos iniciales o de infraestructura, como ocurre con los servicios públicos o los ferrocarriles.

    \item \textbf{Monopolio legal (otorgado por el Estado):} Establecido por ley, concede derechos exclusivos a una empresa o entidad estatal, como los servicios postales, las patentes o industrias estatales como la sal o el tabaco (No confundir con Monopolio estatal).

    \item \textbf{Monopolio de recursos:} Ocurre cuando una empresa o estado controla un recurso crítico, como la sal en la antigua Roma o China, o los diamantes en el caso de De Beers.

    \item \textbf{Monopolio tecnológico:} Surge al poseer una tecnología única o patentes que impiden el ingreso de competidores al mercado (por ejemplo, Microsoft en sus inicios o ciertas patentes farmacéuticas).

    \item \textbf{Monopolio estatal:} El Estado controla directamente la producción y venta de ciertos bienes, como la sal en China o el tabaco en Francia.

    \item \textbf{Monopolio estratégico:} Se logra mediante tácticas empresariales como fusiones, adquisiciones o contratos exclusivos. Ejemplos representativos son Standard Oil y AT\&T.

    \item \textbf{Monopolio colonial o comercial:} Implica la concesión de derechos exclusivos de comercio, usualmente por potencias coloniales, como ocurrió con las Compañías Británica y Holandesa de las Indias Orientales.
\end{itemize}
        \item \textbf{Sistemas complejos, redes, y modelos biológicos de ecosistemas:}
        Después de revisar los antecedentes históricos en la sección pasada, se realizó una investigación sobre Sistemas complejos, redes y modelos biológicos de ecosistemas. El objetivo fue extraer los principales objetos y elementos que podrían usarse en el desarrollo del framework, los cuales se enumeran a continuación:
        \begin{itemize}
        	\item \textbf{Sistemas complejos:}
			\begin{itemize}
		    \item \textbf{Emergencia y no linealidad:} En los sistemas complejos, la emergencia y la no linealidad permiten que los monopolios surjan de interacciones dinámicas entre agentes heterogéneos, cuyos comportamientos agregados no son predecibles desde el nivel individual.

    \item \textbf{Modelado basado en agentes (ABM):} El modelado basado en agentes (ABM) permite simular estrategias de empresas en mercados digitales, mostrando cómo la competencia por recursos y la adaptación a cambios regulatorios inducen concentración.

    \item \textbf{Equilibrio y resiliencia :} Los mercados, al igual que los ecosistemas biológicos, pueden tender a estados de equilibrios inestables o generar monopolios naturales como resultado de efectos de red.
			\end{itemize}
	\item \textbf{Teoría de redes}
		\begin{itemize}
    \item \textbf{Efectos de red y preferencial attachment:} Desde la teoría de redes, los efectos de red y el \textit{preferential attachment} explican cómo plataformas como Amazon o Facebook crecen exponencialmente y crean barreras de entrada.

    \item \textbf{Estructuras scale-free:} La concentración de nodos centrales (empresas dominantes) explica la distribución de poder en mercados digitales

    \item \textbf{Interdependencia: } Las redes de proveedores, consumidores y reguladores forman sistemas interconectados donde cambios locales afectan globalmente el sistema.	
		\end{itemize}
		
	\item \textbf{Modelos biológicos}
		\begin{enumerate}
    \item \textbf{Relaciones ecológicas:} 		
			\begin{itemize}
    \item \textbf{Depredación:} Empresas dominantes pueden eliminar competidores mediante precios agresivos, similar a la exclusión competitiva en ecología
        \item \textbf{Simbiosis y parasitismo:} Alianzas estratégicas (ej. acuerdos de exclusividad) o prácticas como el freeriding reflejan interacciones mutualistas o parasitarias
\end{itemize}			    
    \item \textbf{Nicho ecológico:}Los monopolios ocupan espacios de mercado únicos, controlados mediante innovación o acceso a recursos críticos (ej. datos en la era digital) 		
    \item \textbf{Selección natural: } Las empresas con ventajas adaptativas (tecnología, costos) sobreviven y dominan, mientras otras son desplazadas
		\end{enumerate}

    

    \item \textbf{Ejemplo integrado: } Un modelo de monopolio digital podría combinar:
    \begin{itemize}
    \item \textbf{ABM} (AGENT-BASED MODELS) para simular agentes con estrategias de precios dinámicos (ECONOMÍA COMPUTACIONAL BASADA EN AGENTES)
        \item \textbf{Redes scale-free} ara representar el crecimiento de usuarios en plataformas    
                \item \textbf{Dinámica depredador-presa} para analizar cómo una empresa absorbe competidores    
    \end{itemize}			
\end{itemize}

        
        \item \textbf{Teoría de la antifragilidad en contextos económicos.}
        \begin{itemize}
    \item \textbf{Conceptos esenciales de antifragilidad aplicables a monopolios}
    \begin{enumerate}
        \item Antifragilidad vs. resiliencia: Un sistema antifrágil no solo resiste shocks y volatilidad, sino que se fortalece y mejora ante ellos, a diferencia de un sistema meramente resiliente que solo sobrevive o se recupera.

    \item No linealidad: Las respuestas de los monopolios ante perturbaciones pueden ser no lineales; pequeñas crisis pueden no afectarlos, pero grandes disrupciones pueden hacerlos aún más fuertes, por ejemplo, eliminando competidores débiles o permitiendo la adquisición de innovaciones.

    \item Opcionalidad: Los monopolios antifrágiles mantienen flexibilidad estratégica, diversificando productos o mercados para adaptarse y beneficiarse de la incertidumbre y los cambios regulatorios o tecnológicos.

    \item Gestión de riesgos y tolerancia a la volatilidad: Un monopolio antifrágil aprovecha la volatilidad del mercado, la incertidumbre y los factores estresantes (como crisis económicas o cambios tecnológicos) para innovar, optimizar procesos o consolidar su posición.

    \item Descentralización interna: Aunque los monopolios tienden a la concentración, pueden estructurarse internamente de manera descentralizada (por divisiones, filiales, etc.) para distribuir el riesgo y evitar puntos únicos de falla.

    \item Estrategia de barra (barbell strategy): Consiste en combinar actividades extremadamente seguras con apuestas especulativas, permitiendo a los monopolios protegerse ante grandes shocks y, al mismo tiempo, aprovechar oportunidades de alto riesgo y alta recompensa.

    \item Innovación y aprendizaje de fracasos: Los monopolios antifrágiles fomentan la experimentación y el aprendizaje de errores, utilizando los fracasos como fuente de mejora continua.    
    \end{enumerate}
    \item \textbf{Implicaciones para el modelado de monopolios}
    \begin{itemize}
        \item Simulación de shocks y volatilidad: Es necesario modelar cómo el monopolio responde y se adapta a crisis, cambios regulatorios, disrupciones tecnológicas y ataques competitivos, evaluando si se debilita (frágil), resiste (robusto) o se fortalece (antifrágil).

    \item Análisis de estructura y flexibilidad: Incorporar la capacidad del monopolio para reestructurarse, diversificarse y mantener opciones abiertas frente a la incertidumbre.

    \item Evaluación de estrategias de riesgo: Analizar cómo el monopolio gestiona su exposición a eventos extremos (cisnes negros) y si su estructura le permite beneficiarse de ellos en vez de solo sobrevivir.
    \end{itemize}
\end{itemize}    

    \item Definir operacionalmente conceptos clave: monopolio, emergencia, voluntad centralizada, condiciones sistémicas, antifragilidad.
    \begin{itemize}
    \item \textbf{Referir al texto anexo sobre la implementacion del framework}
    \end{itemize}
    \item Estudiar modelos dinámicos relevantes: ecuaciones de Lotka-Volterra, redes complejas, sistemas autoorganizados.
    \begin{itemize}
    \item \textbf{Referir al texto anexo sobre la implementacion del framework.}
    \end{itemize}
     
    \end{itemize}
\item \textbf{Etapa II: Diseño del Modelo Sistémico} 
	\begin{itemize}
	\item \textbf{Identificar variables clave:}
	oferta, demanda, regulación, escalabilidad, efectos
de red, asimetrías de información, etc. \textbf{Referir al texto anexo sobre la implementacion del framework}
    \item Opciones consideradas para seleccionar el marco modelador más adecuado:
    \begin{itemize}
        \item Sistemas dinámicos no lineales.
        \item Modelos de redes dinámicas.
        \item Modelos basados en agentes (ABM).
        \item \textbf{Modelos híbridos.} (Opción seleccionada) 
    \end{itemize}
	\end{itemize}
Todos los requerimientos de la Etapa II se han satisfecho en el texto anexo del Framework.

\item \textbf{Etapa III: Validación Empírica y Entrenamiento} 
	\begin{itemize}
    \item Recopilar casos históricos y conjuntos de datos reales: Standard Oil, East India Company, Amazon, etc. (anexo datos en repo)
    \item Ajustar parámetros del modelo utilizando técnicas de calibración: optimización bayesiana, simulación Monte Carlo, aprendizaje automático. (referir al jupyter notebook anexo)
    \item Validar el modelo:
    	\begin{itemize}
        \item Comparar resultados simulados con trayectorias históricas.
        \item Evaluar capacidad predictiva y explicativa.
        (referir al jupyter notebook anexo)
    	\end{itemize}
	\end{itemize}
	
\item \textbf{Etapa IV: Evaluación de Antifragilidad} 
	\begin{itemize}
    \item Definir métricas de antifragilidad adaptadas al entorno económico.
    \item Someter el modelo a perturbaciones controladas: choques regulatorios, tecnológicos o financieros.
    \item Clasificar los tipos de monopolio según su respuesta: frágiles, robustos, antifrágiles.
	\end{itemize}   
\end{itemize}



\subsubsection*{. Revisar el Resultado}

\begin{itemize}
    \item Evalúa la solución obtenida y asegúrate de que responde al problema planteado: 
\textbf{Podemos concluir que el modelo aproxima suficientemente bien a las condiciones necesarias para la descripción de "apertura monopólica".} 
\end{itemize}

\section*{Fase 3: Entregables}

Prepara los siguientes entregables documentando todo el proceso:

\subsubsection*{1. Diseño de Prompts y Prompts Usados}


\begin{itemize}
    \item Presentados en el archivo: \textbf{promptsIniciales.pdf}
\end{itemize}


\subsubsection*{2. Escenario del ABP}

\begin{itemize}
    \item Expuesto en el presente archivo.
\end{itemize}

\subsubsection*{3. Documentos de Resolución del ABP}

Siguiendo el esquema de Polya se presentan los siguientes documentos:

\begin{itemize}
    \item abp1
    \item Planeacion
    \item promptsIniciales
    \item bilbiografiaABP1
    \item Framework(incluye resultados obtenidos).
\end{itemize}
Adicionalmente se agregan los links al repositorio con el codigo y datasets usados para validar el modelo.


\subsection*{Conclusión}

A partir del análisis realizado, se concluye que los monopolios pueden tanto ser \textit{creados} deliberadamente como \textit{emergentes}, dependiendo de las condiciones estructurales del sistema económico en cuestión. En efecto, se ha observado que existen monopolios promovidos y sostenidos por estructuras estatales (\textit{state-promoted}), así como otros que surgen de forma autoorganizada (\textit{self-emergent}), impulsados por dinámicas internas del mercado, como la acumulación de ventajas competitivas, economías de escala, control de infraestructura o recursos clave, o el efecto red.

Al modelar los mercados como sistemas dinámicos análogos a ecosistemas biológicos, se identifican patrones recurrentes de comportamiento no lineal, dependencia de condiciones iniciales y fenómenos de retroalimentación que permiten comprender mejor la emergencia de estructuras monopólicas. Esta aproximación sistémica permite también diferenciar entre monopolios \textit{naturales}, que surgen como la solución más eficiente en contextos de altos costos fijos o limitación de recursos, y monopolios \textit{artificiales}, sostenidos por barreras estratégicas o jurídicas impuestas externamente.

Un hallazgo clave de esta investigación es que ciertos monopolios pueden exhibir propiedades de \textit{antifragilidad}: en lugar de simplemente resistir las disrupciones del mercado, se fortalecen activamente ante ellas, ya sea mediante la absorción de competidores debilitados, la integración vertical de cadenas de valor, o la adaptación estratégica a entornos regulatorios y tecnológicos cambiantes. No obstante, esta antifragilidad a nivel empresarial no necesariamente se traduce en resiliencia sistémica. Por el contrario, \textbf{implica una transferencia estructural:} la capacidad de absorber shocks se desplaza del sistema económico al actor dominante, dejando al conjunto del sistema más expuesto a disrupciones inesperadas. Así, el fortalecimiento del monopolio puede ir acompañado de una creciente fragilidad sistémica, al concentrar riesgos, inhibir la competencia y reducir la diversidad adaptativa del mercado. Esta paradoja plantea retos significativos para la regulación económica, ya que un monopolio antifrágil no sólo persiste frente a las crisis, sino que prospera gracias a ellas, a costa de debilitar la capacidad del sistema para regenerarse, adaptarse y evolucionar.

En definitiva, entender el fenómeno de los monopolios requiere ir más allá del marco legal o regulatorio. Es necesario un enfoque que combine herramientas de distintas disciplinas —como la teoría de sistemas, la economía evolutiva y el análisis estructural— para poder reconocer las condiciones que dan lugar a su aparición, permanencia o eventual disolución. Solo así es posible evaluar con claridad si estas estructuras fortalecen o debilitan la capacidad del sistema económico para enfrentar incertidumbres y transformaciones profundas.





\end{document}
